\documentclass[11pt, a4paper]{article}
\usepackage[utf8]{inputenc}
\usepackage[margin=1in]{geometry}
\usepackage{amsmath}
\usepackage{booktabs}
\usepackage{graphicx}
\usepackage{hyperref}

\title{\textbf{War \& the Shadow Economy}}
\author{DERRAG Mariya}
\date{March 2026}

\begin{document}

\maketitle

\begin{abstract}
This research investigates the relationship between armed conflict and the emergence of parallel shadow economies. Using a dataset of 100,000 conflict records (1939--2026), we demonstrate that war systematically redirects economic activity into a shadow system. Through linear regression, we establish that while war itself provides a structural baseline for informal growth, the currency black market gap serves as the primary measurable predictor of its magnitude.
\end{abstract}

\section{Introduction}
Armed conflict is traditionally studied via GDP collapse and hyperinflation. This research approaches war as an \textit{economic transformation} that produces a parallel economy characterized by structural profitability. We seek to identify the factors determining the growth of the informal economy during conflict.

\section{Data \& Methodology}
The analysis utilizes the \textit{War Economic \& Livelihood Impact Dataset} ($N = 100,000$). The key constructed variable is:
\begin{equation}
    \Delta \text{IE} = \text{IE}_{\text{War}} - \text{IE}_{\text{Pre-War}}
\end{equation}
Where $\Delta \text{IE}$ represents the percentage point growth in the informal economy.

\section{Results}
\subsection{Uniformity of Black Market Activity}
The distribution of black market intensity and goods is nearly uniform. Profiteering is documented in 70\% of all conflicts, regardless of region or conflict type.

\subsection{Correlation Analysis}
Traditional economic indicators show near-zero correlation with informal growth:
\begin{itemize}
    \item GDP Change: $r = +0.001$
    \item Inflation Rate: $r = -0.000$
    \item \textbf{Currency Black Market Gap: $r = +0.603$}
\end{itemize}

\section{Linear Regression Analysis}
To quantify the relationship between currency distortions and shadow economy expansion, we employ a Simple Linear Regression (SLR) model.

\subsection{Model Specification}
The growth of the informal economy is modeled as:
\begin{equation}
    \text{IE\_Growth}_i = \beta_0 + \beta_1 (\text{Currency\_Gap}_i) + \epsilon_i
\end{equation}

\subsection{Estimated Regression Equation}
Based on the OLS estimation, the predictive function is:
\begin{equation}
    \hat{y} = 19.39 + 0.0586x
\end{equation}
For every 1\% increase in the currency gap ($x$), the informal economy growth ($\hat{y}$) increases by 0.0586 percentage points. The intercept $\beta_0 = 19.39$ represents the structural baseline growth driven by the existence of war alone.

\subsection{Statistical Diagnostics}
\begin{table}[ht]
\centering
\caption{Regression Coefficients and Significance}
\begin{tabular}{lcccc}
\toprule
\textbf{Variable} & \textbf{Estimate} & \textbf{Std. Error} & \textbf{t-value} & \textbf{Pr($>|t|$)} \\
\midrule
(Intercept) & 19.39 & 0.04927 & 393.5 & $< 2.2 \times 10^{-16}$ \\
Currency Gap & 0.05861 & 0.0002453 & 239.0 & $< 2.2 \times 10^{-16}$ \\
\bottomrule
\end{tabular}
\end{table}

\begin{itemize}
    \item \textbf{Fisher F-Test:} $F = 57,100$ ($p < 2.2 \times 10^{-16}$), indicating a highly significant model.
    \item \textbf{Goodness of Fit:} $R^2 = 0.3635$. The model explains 36.4\% of the variance.
    \item \textbf{Durbin-Watson:} $DW = 2.0008$, confirming no autocorrelation in residuals.
\end{itemize}

\section{Discussion \& Conclusion}
War creates a replacement system immune to formal economic conditions. The currency black market gap is the only quantifiable predictor of this shift, yet the 19.39pp intercept suggests that shadow economies are a structural guarantee of armed conflict.

\end{document}
